\documentclass[11pt]{article}
\usepackage{amsmath,booktabs,hyperref}
\title{Typed Subscription Retirement Is Not Process Retirement}
\date{2026-11-27}
\begin{document}
\maketitle

\paragraph{Audit and scope.}
This is an offline development replay of an evidence dependency in
a simulated cave inspection, not a deployed binary or a rendered
observation. The subject is whether the same fresh supplied finite
point remains finite at the source, after the auxiliary depth bridge
is no longer running. The controlled comparison holds camera
configuration, injected pre-mask depth and XYZ, current point
consumer, and absence of other matching native depth subscribers fixed.
Time, reliability of the notice to the robot, and physical traversability
are not outputs of this test.

\paragraph{Where the dependency arises.}
The pinned SubT CoSTAR launch runs a separate
\texttt{ros\_ign\_image/image\_bridge} for the front depth image, and a
one-way \texttt{parameter\_bridge} for points. The contemporary reference
Melodic image handler advertises into ROS and then unconditionally
requests the Ignition subscription. The downstream optical wrapper
can change the ROS subscription on demand without this handler
withdrawing the Ignition subscription. Process launched, a successful
registration visible to the native publisher, and a callback delivered
to a ROS reader are therefore three different cuts. The checked historical
Dockerfile selects the family of packages but does not identify the
installed binary hashes. In particular a default bridge that has already
registered clips in the no-ROS-recorder arm too; adding an ordinary
ROS subscriber need not create another native depth connection.

\paragraph{What the publisher counts.}
In the selected Transport 9 source, a valid publisher's
\texttt{HasConnections} consults local subscribers and retained remote
entries for the exact topic and message type (also admitting a generic
type). A registration addressed to the publisher process adds an entry;
a named end removes the specified topic/process/node entry. Discovery
timeout and process BYE initiate a callback bearing only the process
identity. In that branch the shared transport deliberately retains its
remote-subscriber table. In contrast, a topic-and-node disconnection
uses the named removal operation. These are conditional effects of
the delivered callbacks, not a guarantee that a real graceful stop
always delivers the named end. Process-level detection by itself does
not remove the subscription from the query used by the sensor.

For a single initially registered matching remote depth subscriber,
no local one, and no other relevant event between the delivered signal
and the next fresh RGB-D update, the result is:
\[
\begin{array}{lcc}
\text{delivered event} & \text{query} & \text{37 m point fields}\\
\hline
\text{named end to owning publisher} & \text{false} & \text{finite XYZ}\\
\text{process-only expiry or BYE} & \text{true} & +\infty,+\infty,+\infty
\end{array}
\]
A new process can register and remove its own named entry without
removing the old one. An additional matching entry also prevents
named retirement of the first from unblocking the point. A different
topic or concrete type does not suffice. The ordinary reference render
and depth-image planes remain 300 and 10 m respectively; both arms
preserve the supplied 5 m and exactly 10 m controls.

\paragraph{Executed discriminator.}
The test extracts the Handler, the exact callback payload construction
for expiration and process BYE, the publisher's query and three
subscriber cuts, and the original templated topic storage with a minimal
transport/type shim. It feeds the query to the pinned camera clip branch
and runs the prior native copy, mask and packed XYZ routines. An
independent extensional ledger matches the cut sequence at each fresh
frame, separately from the packed value checks. Five termination arms,
five extra-entry arms, three replacement arms, two initial-report arms
and two local-connection arms yield 300 histories, 1200 frames and
10800 packed XYZ field checks. Fixed-address entries and no concurrent
events during a frame are test constraints. The branch selects no
arbitrary network input; no discovery transport stack is launched.

\paragraph{Use at the decision cut.}
Neither bridge liveness nor its projected ROS subscription count
qualifies which clip was applied to a particular frame. The sufficient
local discriminator for this mask law is the actual selected clip
branch tied to that frame, together with camera configuration and
frame provenance, not a poll of the same branch much later. In a live
integration that discriminator still has to be obtained, timed,
and transported through the decision cut under the same faults.
An infinite value is not evidence of open passage in either case:
in the clipped arm the same exported XYZ fits a finite distant return
and an already-infinite renderer sample. A finite positive return may
merit a contact test but is not swept-body clearance. Replacing the
subscriber-membership guard with a policy fixed to the declared
point contract would remove this particular cross-effect, but is not
stock and has not been pursued as a simulator patch. The independent
cave benchmark and contract acceptance remain open.
\end{document}
